\section{Constructive paired packets and reciprocal adjacent divisors}
\label{sec-pa18-paired-packets}

All factors in this section belong to the complete actual squarefree
$Q=\operatorname{rad}(c-1)$. For $c=2^e3^f$, retain (NG0),
$A=2^{e-1}$, $B=3^{f-1}$, $D=\log(AB/Q)>0$, and the same original
optimum $D+\Gamma$. The elementary divisor assertions also apply to
any squarefree $Q>1$ and real $A,B\geq1$ with $AB>Q$.

\paragraph{PA1: a certificate with at most $m$ exact updates.}
Suppose there is an exact integer factorization
\[
 Q=R\prod_{i=1}^mU_iV_i,\qquad m\geq1,\quad
 R\geq1,\quad 1<U_i<V_i.
\]
All its nonunit factors are squarefree and pairwise coprime. They
may be composite packets. Choose an actual positive divisor $K\mid R$
and define
\[
 h_{\min}=K\prod_iU_i,\qquad h_{\max}=K\prod_iV_i,
 \qquad \rho=\max_i(V_i/U_i)>1.
\]
If
\[
 h_{\min}\leq A,\qquad h_{\max}\geq Q/B,              \tag{PA1}
\]
then
\[
 \Gamma\leq\frac{\max(0,\log\rho-D)}2,\qquad
 D+\Gamma\leq\max\left(D,\frac{D+\log\rho}{2}\right).
                                                               \tag{PA2}
\]
In particular $\rho\leq AB/Q$ implies $\Gamma=0$.

\emph{Proof.} For each subset $J$ of packet indices, the two integers
\[
 H_J=K\prod_{i\notin J}U_i\prod_{i\in J}V_i,\qquad
 Q/H_J=(R/K)\prod_{i\notin J}V_i\prod_{i\in J}U_i
\]
recover $Q$ exactly. Every prime is used in exactly one output, so
this is a legal subfamily of the original divisor certificate.
Only $m+1$ members are necessary. Fix the given order and put
\[
 H_j=K\prod_{i\leq j}V_i\prod_{i>j}U_i,
             \qquad 0\leq j\leq m.
\]
This chain is strictly increasing from $h_{\min}$ to $h_{\max}$,
with $H_j/H_{j-1}=V_j/U_j\leq\rho$. Let $j$ be the first index
with $H_j\geq Q/B$, which exists by (PA1). If $H_j\leq A$, a
central divisor has been found. Otherwise $j\geq1$, since
$H_0\leq A$, and
\[
 H_{j-1}<Q/B\leq A<H_j.
\]
The distances from the central interval to these two genuine
divisors sum to $\log(H_j/H_{j-1})-D\leq\log\rho-D$.
Their minimum bounds the optimum over all actual divisors, proving
the first assertion of (PA2). Adding $D$ proves the second. $\square$

After this factor certificate is supplied, the construction is an
exact bounded algorithm. Start at $H_0$ and update
$H_j=(H_{j-1}/U_j)V_j$; division is exact because that packet has
not yet been replaced. At most $m$ updates find a central divisor
or its two brackets. For integer endpoints, all tests are integer
comparisons $BH_j\geq Q$ and $H_j\leq A$. In the bracket case
compare the positive rational numbers $Q/(BH_{j-1})$ and $H_j/A$
to choose the smaller logarithmic penalty. There is no $2^m$ search
or floating-point comparison. This is not an efficient factorization
algorithm or a universal theorem supplying suitable packets.

\paragraph{PA2: an actual instance and its conditional scope.}
For the completely certified $c=2^{41}3^{26}$ of NG3, choose
\[
 \begin{split}
 U_1&=2729,\qquad V_1=3073=7\cdot439,\\
 R&=857\cdot292183\cdot380261663,\qquad K=380261663.
 \end{split}
\]
These factors recover the complete $Q$. The exact inequalities
\[
 K\cdot2729\leq2^{40},\qquad K\cdot3073\cdot3^{25}\geq Q,
 \qquad 3073Q<2729\cdot2^{40}\cdot3^{25}
\]
give (PA1) and $\rho<AB/Q$, hence $\Gamma=0$. The endpoint is
nonsquare, while its global gap parameter $439/7$ exceeds $AB/Q$.
The local packet certificate is therefore a strictly sharper
sufficient test in this actual example. Its added information is a
pair of close composite packets in the required location.

For a sequence of actual endpoints with $Q\to\infty$, if
$D/\log Q\to0$ and certificates satisfying (PA1) can be supplied
with $\log\rho/\log Q\to0$, then (PA2) gives
$(D+\Gamma)/\log Q\to0$. Likewise the two inequalities
\[
 D\leq\epsilon\log Q+C,\qquad
 \log\rho\leq\epsilon\log Q+C
\]
give $D+\Gamma\leq\epsilon\log Q+C$ with the same constants.
Neither hypothesis is established for all endpoints. In this
stratum
\[
 D=\log c-\log\operatorname{rad}(c(c-1)),
\]
so the first hypothesis already contains the restricted scalar ABC
problem. The paired-packet supply is a separate unproved arithmetic
condition, not a consequence of one example or a moderate largest
prime factor.

\paragraph{PA3: the general reciprocal obstruction.}
Let $Q=qR$ be squarefree with $q$ prime, let $d>1$ be a divisor of
$R$, and set
\[
 t=\max\{u:u\mid R,\ u<d\}.
\]
This maximum exists since $1<d$. If
\[
 dtq\leq R<d^2q,                                    \tag{PA3}
\]
then $R/d,dq$ are consecutive actual divisors of $Q$. If also
$dq>\max(A,B)$, then
\[
 \Gamma=\log\frac{dq}{\max(A,B)},\qquad
 D+\Gamma=\log\frac{\min(A,B)}{R/d}.                 \tag{PA4}
\]

\emph{Proof.} If a divisor $H$ contains $q$ and is below $dq$,
write $H=qv$ with $v\mid R,v<d$; then $H\leq qt\leq R/d$.
If it omits $q$ and is above $R/d$, then $R/H<d$, hence
$R/H\leq t$ and $H\geq R/t\geq dq$. Both endpoints divide $Q$,
are strictly ordered, and have product $Q$. The central interval
and its distances are therefore exactly as in NG4, proving (PA4).
$\square$

There is a prime-prefix specialization without a search for $t$.
If $d$ is a prime factor of $R$ and $T$ is the product of every
prime of $R$ smaller than $d$, the additional condition $T<d$
gives $t=T$. Indeed, a divisor below $d$ cannot contain a larger
prime, and $T$ itself is such a divisor. Thus
$dTq\leq R<d^2q$ suffices; NG4 is the case that $d$ is the second
smallest prime of $R$.

This invariant uses the full actual $R=Q/q$ and a certified divisor
threshold. Largest-prime size alone does not contain that information.
An exact obstruction of this kind prevents a PA1 certificate with
$\rho\leq AB/Q$ for the same central interval, since their proved
conclusions would contradict one another. It does not exclude a
different certificate with different proved rules, nor the parent
ABC routes. All accounting above uses actual radicals and whole
prime packets. No infinite fixed-$\epsilon$ counterexample or new
Lean theorem is claimed.
